\documentclass[12pt]{article}
\usepackage{ceai}

\begin{document}

\title{CE 488/5588 Homework 04\\[2in]
\textbf{Name: \rule{2in}{0.15mm}}}
\author{}
\date{}
\maketitle

\newpage
\doublespacing

% ============================================================
\section*{Problem 1: Perceptron Model and Decision Geometry}

Consider a perceptron model for binary classification.
The input is a two-dimensional feature vector
\[
  \bm{x} =
  \begin{bmatrix}
    x_1 \\
    x_2
  \end{bmatrix}.
\]

After training, the learned parameters are
\[
  [w_1, w_2, b]^T = [-4, -3, 10]^T.
\]

Assume the binary class labels are \( c \in \{0,1\} \).
When using the Heaviside activation, define
\[
  H(u) =
  \begin{cases}
    1, & u \ge 0, \\
    0, & u < 0.
  \end{cases}
\]

\begin{enumerate}
  \item Write down the linear decision function \( u(\bm{x}) \).

  \item Derive the equation of the decision boundary \( u(\bm{x}) = 0 \).
    Express the boundary in the form \( x_2 = f(x_1) \).
    Sketch the line in the \( x_1\text{-}x_2 \) plane and clearly indicate the two half-spaces.

  \item Using the Heaviside activation defined above, write the full perceptron model \( m(\bm{x}) \).

  \item For the following two data points:
    \[
      \text{A: } [1,4]^T,
      \qquad
      \text{B: } [-2,2]^T,
    \]
    compute \( u(\bm{x}) \) and determine the predicted class label for each point.

  \item Now suppose the activation function is replaced by a ReLU:
    \[
      \mathrm{ReLU}(u) = \max(0,u).
    \]
    For each of the two points above, compute the scalar output
    \[
      m(\bm{x}) = \max(0, u(\bm{x})).
    \]
    (Report the numerical output value; do not convert it to a class label.)
\end{enumerate}

\newpage
% ============================================================
\section*{Problem 2: Multilayer Perceptron Architecture and Parameter Counting}

A civil engineer constructs a supervised learning exercise with the following setup:

\begin{itemize}
  \item Each input data point is a \( 3 \times 1 \) vector.
  \item The output is binary classification with labels \( \{0,1\} \).
  \item The dataset contains 1000 labeled data points.
  \item The neural network is a fully connected multilayer perceptron (MLP).
\end{itemize}

The network architecture is defined as:

\begin{itemize}
  \item Input layer: dimension 3.
  \item Hidden layer 1: 5 neurons.
  \item Hidden layer 2: 5 neurons.
  \item Output layer: 1 neuron (binary logit).
\end{itemize}

Bias parameters are included at every non-input layer.

\begin{enumerate}
  \item Determine the total number of trainable parameters in this MLP, including all weights and biases. Clearly show how you compute the count layer by layer.

  \item Draw a clear diagram of the network architecture, indicating:
    \begin{itemize}
      \item Number of neurons in each layer,
      \item Fully connected structure between layers,
      \item Bias connections.
    \end{itemize}
\end{enumerate}

\newpage
% ============================================================
\section*{Problem 3: Perceptron vs. Linear SVM}

Recall that both a perceptron (with hard sign activation) and a linear Support Vector Machine (SVM) produce a decision function of the form

\[
  m(\bm{x}) = \mathrm{sign}\bigl(u(\bm{x})\bigr),
  \qquad
  u(\bm{x}) = \langle \bm{w}, \bm{x} \rangle + w_0.
\]

Despite sharing the same functional form at inference time, their learning principles differ.

List and briefly explain several key differences between a perceptron and a linear SVM in the binary classification setting.
You may consider discussing:

\begin{itemize}
  \item The learning objective,
  \item Margin definition and margin width,
  \item Loss functions,
  \item Behavior under non-separable data,
  \item Optimization properties.
\end{itemize}

Provide concise but conceptually clear explanations.

\newpage
% ============================================================
\section*{Problem 4: Polynomial Kernel and Feature Lifting}

\textbf{Required for Graduate Students; and self-learning ot Topic07 - Section 4 maybe needed for a mathematical exploration.}

Consider a polynomial kernel SVM defined by
\begin{align*}
  K(\bm{x}, \bm{x}')
  &= \langle \Phi(\bm{x}), \Phi(\bm{x}') \rangle \\
  &= \bigl(\langle \bm{x}, \bm{x}' \rangle + 1 \bigr)^d,
\end{align*}
where \( d \) is the degree of the polynomial.

If the original input space has dimension \( D \), then the transformed feature space has dimension

\[
  D_k = \binom{D + d}{d}.
\]

\begin{enumerate}
  \item Suppose \( D = 2 \) and \( d = 2 \).
    Compute \( D_k \).
    What is the dimension of the transformed feature space?

  \item For the specific case
    \[
      \bm{x} = [x_1, x_2]^T,
      \qquad
      K(\bm{x}, \bm{x}') = (\langle \bm{x}, \bm{x}' \rangle + 1)^2,
    \]
    derive an explicit feature mapping \( \Phi(\bm{x}) \).

    Your answer must satisfy
    \[
      \Phi(\bm{x})^T \Phi(\bm{x}')
      =
      (\langle \bm{x}, \bm{x}' \rangle + 1)^2.
    \]

    Show your expansion and verify that your mapping is correct.
    (Any valid equivalent feature map that satisfies the identity is acceptable.)

    You may consult external tools for brainstorming, but you must verify all mathematical steps \footnote{ For example, you may prompt in Gemini or Chatgpt:
      \begin{quote}
        ``In a polynomial kernel SVM, if the input feature is $\bm{x} = [x_1, x_2]^T$, and the kernel is $k(\bm{x}, \bm{x}') = (\langle \bm{x}, \bm{x}'\rangle + 1)^2$, i.e., using a 2nd-order polynomial kernel, what does the transformation mapping look like?''
    \end{quote}}
\end{enumerate}

\end{document}
